% ==============================================================================
% Fichier : front/03_ndap_sukulu.tex
% Titre   : NDAP SUKULU, l'environnement compagnon du cours
% Auteur  : MINKA MI NGUIDJOÏ Thierry Emmanuel
% Version : 1.0 -- novembre 2024
% ==============================================================================
% Texte commun aux cours de l'auteur, mis à jour en novembre 2024 (version en
% ligne, banques de 100 questions, compte de démonstration) ; le tableau des
% labs est propre au cours (front/03b_labs.tex).
% ==============================================================================

\chapter*{NDAP SUKULU, la salle de classe compagnon}
\addcontentsline{toc}{chapter}{NDAP SUKULU, la salle de classe compagnon}
\markboth{NDAP SUKULU}{NDAP SUKULU}

\begin{center}
    \includegraphics[width=0.62\textwidth]{figures/logos/logo_ndap_sukulu.png}
\end{center}

\section*{Un nom venu du pays bassa}

\begin{quote}\itshape
Le soir tombe sur Éséka. Un enfant revient des champs avec son grand-père et
lui demande pourquoi il doit marcher si loin, chaque matin, pour aller à
l'école. Le vieil homme s'arrête au bord du chemin et répond, en bassa~:
« \textup{Ndap sukulu}, la salle de classe, ce ne sont pas les murs. Ce sont
les murs qui viennent après. La salle de classe, c'est l'endroit où quelqu'un
t'attend pour t'apprendre ce qu'il sait. » L'enfant insiste~: « Et qui
m'attend~? » Le grand-père sourit~: « \textup{Malet yon}. Ton enseignant. »
\end{quote}

En bassa, la langue paternelle de l'auteur, \textit{ndap sukulu} désigne la
salle de classe et \textit{malet yon} signifie « ton enseignant ».
L'environnement compagnon de ce cours porte ce nom parce qu'il prolonge la
salle de classe au-delà des heures de cours~: il attend l'étudiant le soir,
le week-end, à distance, avec des questions et des pièces à examiner. Il
est édité par MaletYon et mis à la disposition des étudiants et des
praticiens qui souhaitent se perfectionner
(\href{mailto:maletyon@proton.me}{maletyon@proton.me}).

\section*{Ce que l'on y trouve}

\begin{itemize}
    \item Une fiche de questions par chapitre. Chaque tentative tire vingt
          questions d'une banque de cent (QCM, vrai ou faux, textes à
          trous) et dure trente minutes.
    \item Dix laboratoires. Chaque lab place l'étudiant devant les pièces
          d'une affaire du cours (images, captures, journaux, tableurs de
          balayage) et se termine par des points de contrôle.
    \item Un espace enseignant, qui corrige les remises, établit la
          synthèse de la classe et publie les résultats.
\end{itemize}

Les fiches et les labs sont des travaux à faire hors des heures de cours.
Ils ne remplacent ni le livre ni les exercices de fin de chapitre~: ils en
vérifient la maîtrise.

\section*{Comment y accéder}

\begin{enumerate}
    \item En ligne, à l'adresse
          \url{https://maletyon.github.io/cours-fav/ndap-sukulu/}, dans un
          navigateur récent~; sur téléphone, ajouter \texttt{telephone/} à
          la fin de l'adresse.
    \item Hors ligne, en téléchargeant l'archive de l'édition ordinateur
          depuis le dépôt \texttt{MaletYon/cours-fav}, dossier
          \texttt{ndap-sukulu}.
    \item Se connecter avec le matricule et le mot de passe remis par
          l'enseignant. Pour essayer avant d'être inscrit, un compte de
          démonstration est ouvert à tous~: matricule \texttt{ETU001}, mot
          de passe \texttt{eTU001}~; ses essais ne sont jamais transmis à
          l'enseignant.
    \item Pour un lab, vérifier d'abord l'empreinte SHA-256 des pièces
          (\texttt{sha256sum -c MANIFESTE\_SHA256.txt}), puis travailler sur
          une copie~: on n'examine jamais l'original.
\end{enumerate}

\begin{boxinfo}
Les fiches de questions sont aussi fournies au format texte dans le dossier
\texttt{fiches} du dépôt du cours, prêtes à être distribuées. Leurs
corrections figurent dans le livret de correction, à la suite des corrigés
des exercices de chaque chapitre.
\end{boxinfo}

\section*{Les laboratoires du cours}

% Tableau des labs du cours FAV dans NDAP SUKULU (propre au cours).
Le cours compte dix laboratoires, disponibles sous le code \codeUE. Les six
premiers reprennent les pièces de l'enquête NYANGA, les trois suivants le
programme de SOLITRANS, et le dernier conduit une enquête complète sur la
compromission du portail de SOLITRANS en février 2026, de la première trace
à la faille.

\begin{tableaucours}{lL{6.2cm}L{1.6cm}L{2.2cm}}{%
    \tableauHeader{Lab} &
    \tableauHeader{Sujet} &
    \tableauHeader{Chapitre} &
    \tableauHeader{Opération}
}
FAV-LAB-01 & Chaîne de possession et scellés & 2 & NYANGA \\
FAV-LAB-02 & Duplication et hachage & 3 & NYANGA \\
FAV-LAB-03 & Disque~: système de fichiers et artefacts & 4 & NYANGA \\
FAV-LAB-04 & Mémoire~: processus et connexions & 5 & NYANGA \\
FAV-LAB-05 & Réseau~: l'exfiltration & 6 & NYANGA \\
FAV-LAB-06 & Chronologie unifiée & 7 & NYANGA \\
FAV-LAB-07 & Lire un résultat de balayage & 10 & SOLITRANS \\
FAV-LAB-08 & Coter~: CVSS v4.0, EPSS, KEV & 9 & SOLITRANS \\
FAV-LAB-09 & Prioriser et suivre la remédiation & 11 & SOLITRANS \\
FAV-LAB-10 & Enquête intégratrice & 12, 14 & SOLITRANS \\
\end{tableaucours}
\vspace{-0.8em}
\begin{center}{\small\itshape Les laboratoires du cours dans NDAP SUKULU}\end{center}

